% !PRINTED 294
\weberpage{294}
%<<CONT
\[
\begin{aligned}
A_\alpha(\varphi,\psi) &\quad -1,\\
A_\varepsilon(\varphi,\psi) &\quad +1,\\
A'(\varphi,\psi) &\quad -1,\\
E'(\varphi,\psi) &\quad +1.
\end{aligned}
\]

Twice the value of the characteristic is then, if for the sake of simplicity we omit the designation $(\varphi,\psi)$ (§.\,95),
\weq{3}{-A_\alpha + A_\varepsilon - A' + E' = 2k;}
adding to this the equation (2) yields
\weq{4}{k = A_\varepsilon,}
i.\,e.\ \gesperrt{the characteristic is equal to the number of the intersection points of $\varphi, \psi$ lying in the interior of $f$.}

With this we have thus found a means of determining, from the characteristic, the number of roots of the equation
\[
F(z) = 0,
\]
whose geometric images lie in the interior of an arbitrarily given curve. The outer connections added at will play no further part here at all, and can be suppressed entirely.

\begin{center}
§. 97.

Determination of the Characteristic.
\end{center}

To make the theorem proved in the foregoing applicable, we must still show how we can actually determine the characteristics. For this purpose, let us imagine the curve $f$ divided into segments in such a way that in any one of them there lies only one intersection point with the curve $\varphi$ or one intersection point with the curve $\psi$; and let us denote by $a$ (Fig. 13) the beginning, by $b$ the end of such a segment, and by $\xi$ the intersection point of $f$ with $\varphi$ lying in it. Then the point $\xi$ will be a point $E(f;\varphi,\psi)$ if the product $\varphi\psi$ is positive at $a$, and a point $A(f;\varphi,\psi)$ if $\varphi\psi$ is negative at $a$; and with this the

\weberfigureimage{weber-b1-fig13.png}{0.42}
%CONT>>
